\section{Conclusão}

A solução \ac{FAR} da Tevel atende ao requisito de aplicação desenvolvida e testada em ambiente agrícola, com colheita integrada documentada em pomares. O caso reúne sensoriamento, processamento, atuação e gestão digital em uma arquitetura \ac{IoT} voltada à escassez de mão de obra na janela de colheita. Robôs aéreos e suporte terrestre operam com processamento intensivo local e serviços de supervisão e manutenção.

A pesquisa identificou evidências de colheita seletiva e aperfeiçoamento da infraestrutura computacional. Os benefícios potenciais incluem maior capacidade de atender à demanda de colheita e melhor conhecimento da produção. Contudo, as fontes consultadas, predominantemente vinculadas ao fabricante, clientes e parceiros, não permitem quantificar independentemente produtividade, danos, economia, disponibilidade ou impacto energético, limitando conclusões sobre superioridade econômica e agronômica.

Para a engenharia de \acp{VANT}, o caso evidencia a integração entre percepção, controle, energia, comunicação e logística, conectada pela \ac{IoT} a serviços e registros agrícolas. O desempenho deve ser avaliado pela quantidade de fruta comercializável entregue no prazo e pelo custo de manter o sistema operante.